\section*{Аналитический расчёт ОИСКО для прямоугольного ядра}

Пусть
\[
a=\sqrt{3},\qquad
f(x)=K(x)=\frac{1}{2a}\mathbf{1}_{\{|x|\le a\}}.
\]
Здесь $f$ является плотностью равномерного распределения на $[-a,a]$, а
$K$ берётся тем же прямоугольным ядром. Параметр сглаживания будем сразу
обозначать через $\xi$, как в учебнике.

Рассматривается ядерная оценка плотности
\[
\widehat f_\xi(x)=\frac{1}{n\xi}\sum_{i=1}^n
K\left(\frac{x-X_i}{\xi}\right),
\qquad X_i\sim U[-a,a].
\]
Введём масштабированное ядро
\[
K_\xi(x)=\frac{1}{\xi}K\left(\frac{x}{\xi}\right)
=\frac{1}{2a\xi}\mathbf{1}_{\{|x|\le a\xi\}}.
\]
Тогда одно слагаемое в оценке имеет вид $K_\xi(x-X_i)$, а математическое
ожидание оценки равно свёртке:
\[
\mathbb E\widehat f_\xi(x)=(K_\xi*f)(x).
\]

ОИСКО определяется как математическое ожидание интегрального квадрата
отклонения оценки от истинной плотности:
\[
\bar{\delta}_n(\xi)
=\mathbb E\int_{-\infty}^{\infty}
\bigl(\widehat f_\xi(x)-f(x)\bigr)^2\,dx.
\]
Для каждого фиксированного $x$ используем тождество
\[
\mathbb E(\widehat f_\xi(x)-f(x))^2
=\operatorname{Var}\widehat f_\xi(x)
+\bigl(\mathbb E\widehat f_\xi(x)-f(x)\bigr)^2.
\]
После интегрирования по $x$ получаем разложение ОИСКО на дисперсионную
часть и квадрат смещения:
\[
\bar{\delta}_n(\xi)
=\int_{-\infty}^{\infty}\operatorname{Var}\widehat f_\xi(x)\,dx
+\int_{-\infty}^{\infty}
\bigl((K_\xi*f)(x)-f(x)\bigr)^2\,dx.
\]

\section*{Дисперсионная часть}

Для одного наблюдения положим
\[
Y_x=K_\xi(x-X).
\]
Так как наблюдения независимы и одинаково распределены,
\[
\operatorname{Var}\widehat f_\xi(x)
=\frac{1}{n}\operatorname{Var}Y_x.
\]
Следовательно,
\[
\int \operatorname{Var}\widehat f_\xi(x)\,dx
=\frac{1}{n}
\left(
\int \mathbb E Y_x^2\,dx
-\int \bigl(\mathbb E Y_x\bigr)^2\,dx
\right).
\]

Первый интеграл упрощается заменой переменной $u=x-t$ и теоремой Фубини:
\[
\begin{aligned}
\int \mathbb E Y_x^2\,dx
&=\int f(t)\int K_\xi^2(x-t)\,dx\,dt\\
&=\int K_\xi^2(u)\,du
=\frac{1}{\xi}\int K^2(u)\,du.
\end{aligned}
\]
Для прямоугольного ядра
\[
R(K)=\int K^2(u)\,du
=\int_{-a}^{a}\frac{1}{4a^2}\,du
=\frac{1}{2a}.
\]
Кроме того,
\[
\mathbb E Y_x=(K_\xi*f)(x).
\]
Поэтому дисперсионная часть равна
\[
\int \operatorname{Var}\widehat f_\xi(x)\,dx
=\frac{1}{n}
\left(
\frac{R(K)}{\xi}-\|K_\xi*f\|_2^2
\right).
\]

\section*{Свёртка прямоугольных плотностей}

Найдём $(K_\xi*f)(x)$. Оба множителя под интегралом являются постоянными на
своих носителях, поэтому свёртка равна произведению высот на длину
пересечения отрезков
\[
[x-a\xi,x+a\xi]\quad\text{и}\quad[-a,a].
\]
То есть
\[
(K_\xi*f)(x)
=\frac{1}{4a^2\xi}\,\ell(x),
\qquad
\ell(x)=
\operatorname{length}\left([x-a\xi,x+a\xi]\cap[-a,a]\right).
\]
Вид свёртки меняется в зависимости от того, какой из отрезков длиннее.
Если $0<\xi\le 1$, сглаживающее окно короче носителя $f$, поэтому в центре
свёртка совпадает с $f$. Если $\xi\ge 1$, появляется центральная площадка
высоты $1/(2a\xi)$.

При $0<\xi\le 1$
\[
(K_\xi*f)(x)=
\begin{cases}
\dfrac{1}{2a}, & |x|\le a(1-\xi),\\[6pt]
\dfrac{a(1+\xi)-|x|}{4a^2\xi},
& a(1-\xi)<|x|\le a(1+\xi),\\[6pt]
0, & |x|>a(1+\xi).
\end{cases}
\]
Используя чётность функции, получаем
\[
\begin{aligned}
\|K_\xi*f\|_2^2
&=2\int_0^{a(1-\xi)}\frac{1}{4a^2}\,dx
+2\int_{a(1-\xi)}^{a(1+\xi)}
\left(\frac{a(1+\xi)-x}{4a^2\xi}\right)^2 dx\\
&=\frac{1-\xi}{2a}
+\frac{1}{8a^4\xi^2}
\int_{a(1-\xi)}^{a(1+\xi)}
\bigl(a(1+\xi)-x\bigr)^2 dx\\
&=\frac{1-\xi}{2a}
+\frac{1}{8a^4\xi^2}\cdot\frac{(2a\xi)^3}{3}
=\frac{3-\xi}{6a}.
\end{aligned}
\]

При $\xi\ge 1$
\[
(K_\xi*f)(x)=
\begin{cases}
\dfrac{1}{2a\xi}, & |x|\le a(\xi-1),\\[6pt]
\dfrac{a(1+\xi)-|x|}{4a^2\xi},
& a(\xi-1)<|x|\le a(1+\xi),\\[6pt]
0, & |x|>a(1+\xi).
\end{cases}
\]
Аналогично,
\[
\begin{aligned}
\|K_\xi*f\|_2^2
&=2\int_0^{a(\xi-1)}\frac{1}{4a^2\xi^2}\,dx
+2\int_{a(\xi-1)}^{a(1+\xi)}
\left(\frac{a(1+\xi)-x}{4a^2\xi}\right)^2 dx\\
&=\frac{\xi-1}{2a\xi^2}
+\frac{1}{8a^4\xi^2}\cdot\frac{(2a)^3}{3}
=\frac{3\xi-1}{6a\xi^2}.
\end{aligned}
\]
Итак,
\[
\|K_\xi*f\|_2^2=
\begin{cases}
\dfrac{3-\xi}{6a}, & 0<\xi\le 1,\\[6pt]
\dfrac{3\xi-1}{6a\xi^2}, & \xi\ge 1.
\end{cases}
\]

\section*{Квадрат смещения}

Теперь вычислим
\[
\|K_\xi*f-f\|_2^2.
\]
Здесь приходится рассматривать три случая, потому что взаимное положение
центральной площадки свёртки и носителя $f$ меняется при $\xi=1$ и
$\xi=2$.

При $0<\xi\le 1$ в центре отрезка $[-a(1-\xi),a(1-\xi)]$ свёртка совпадает
с $f$. Ошибка возникает только в двух внутренних граничных слоях и двух
внешних хвостах:
\[
\begin{aligned}
\|K_\xi*f-f\|_2^2
&=2\int_{a(1-\xi)}^{a}
\left(
\frac{a(1+\xi)-x}{4a^2\xi}-\frac{1}{2a}
\right)^2 dx\\
&\quad
+2\int_{a}^{a(1+\xi)}
\left(\frac{a(1+\xi)-x}{4a^2\xi}\right)^2 dx.
\end{aligned}
\]
В обоих интегралах после замены переменной получается один и тот же
интеграл по отрезку длины $a\xi$:
\[
\|K_\xi*f-f\|_2^2
=2\int_0^{a\xi}\left(\frac{u}{4a^2\xi}\right)^2du
+2\int_0^{a\xi}\left(\frac{u}{4a^2\xi}\right)^2du
=\frac{\xi}{12a}.
\]

При $1<\xi\le 2$ центральная площадка свёртки
$|x|\le a(\xi-1)$ лежит внутри носителя $f$. Поэтому квадрат смещения
состоит из вклада центральной площадки, наклонной части внутри носителя и
внешнего хвоста:
\[
\begin{aligned}
\|K_\xi*f-f\|_2^2
&=2\int_0^{a(\xi-1)}
\left(\frac{1}{2a\xi}-\frac{1}{2a}\right)^2dx\\
&\quad
+2\int_{a(\xi-1)}^{a}
\left(
\frac{a(1+\xi)-x}{4a^2\xi}-\frac{1}{2a}
\right)^2dx\\
&\quad
+2\int_a^{a(1+\xi)}
\left(\frac{a(1+\xi)-x}{4a^2\xi}\right)^2dx.
\end{aligned}
\]
После вычисления элементарных интегралов:
\[
\|K_\xi*f-f\|_2^2
=\frac{3\xi^3-6\xi^2+6\xi-2}{12a\xi^2},
\qquad 1<\xi\le 2.
\]

При $\xi\ge 2$ весь носитель $f$ лежит внутри центральной площадки
свёртки, а вне носителя $f$ остаются ещё постоянные участки свёртки и
наклонные хвосты. Поэтому
\[
\begin{aligned}
\|K_\xi*f-f\|_2^2
&=2\int_0^a
\left(\frac{1}{2a\xi}-\frac{1}{2a}\right)^2dx\\
&\quad
+2\int_a^{a(\xi-1)}
\left(\frac{1}{2a\xi}\right)^2dx\\
&\quad
+2\int_{a(\xi-1)}^{a(1+\xi)}
\left(\frac{a(1+\xi)-x}{4a^2\xi}\right)^2dx\\
&=\frac{3\xi^2-3\xi-1}{6a\xi^2}.
\end{aligned}
\]
Следовательно,
\[
\|K_\xi*f-f\|_2^2=
\begin{cases}
\dfrac{\xi}{12a}, & 0<\xi\le 1,\\[6pt]
\dfrac{3\xi^3-6\xi^2+6\xi-2}{12a\xi^2}, & 1<\xi\le 2,\\[6pt]
\dfrac{3\xi^2-3\xi-1}{6a\xi^2}, & \xi\ge 2.
\end{cases}
\]

\section*{Итоговая формула для ОИСКО}

Собираем дисперсионную часть и квадрат смещения. Для $0<\xi\le 1$
\[
\frac{R(K)}{\xi}-\|K_\xi*f\|_2^2
=\frac{1}{2a\xi}-\frac{3-\xi}{6a}.
\]
Для $\xi\ge 1$
\[
\frac{R(K)}{\xi}-\|K_\xi*f\|_2^2
=\frac{1}{2a\xi}-\frac{3\xi-1}{6a\xi^2}
=\frac{1}{6a\xi^2}.
\]
Значит,
\[
\bar{\delta}_n(\xi)=
\begin{cases}
\dfrac{\xi}{12a}
+
\dfrac{1}{n}\left(\dfrac{1}{2a\xi}-\dfrac{3-\xi}{6a}\right),
& 0<\xi\le 1,\\[10pt]
\dfrac{3\xi^3-6\xi^2+6\xi-2}{12a\xi^2}
+
\dfrac{1}{n}\cdot\dfrac{1}{6a\xi^2},
& 1<\xi\le 2,\\[10pt]
\dfrac{3\xi^2-3\xi-1}{6a\xi^2}
+
\dfrac{1}{n}\cdot\dfrac{1}{6a\xi^2},
& \xi\ge 2,
\end{cases}
\qquad a=\sqrt{3}.
\]

Проверим согласованность формулы. В точке $\xi=1$ первая и вторая ветви
дают
\[
\frac{1}{12a}+\frac{1}{6an}.
\]
В точке $\xi=2$ вторая и третья ветви дают
\[
\frac{5}{24a}+\frac{1}{24an}.
\]
Значит, $\bar{\delta}_n(\xi)$ непрерывна при $\xi=1$ и $\xi=2$.
Кроме того,
\[
\xi\to 0+:\qquad
\bar{\delta}_n(\xi)\sim \frac{1}{2an\xi}\to\infty,
\]
то есть слишком малое окно даёт большой дисперсионный вклад. При
$\xi\to\infty$
\[
\bar{\delta}_n(\xi)\to\frac{1}{2a}
=\int_{-\infty}^{\infty}f^2(x)\,dx
=\frac{1}{2\sqrt{3}},
\]
что соответствует чрезмерному сглаживанию: оценка становится слишком
широкой и плохо приближает исходную плотность.

\section*{Оптимальный параметр и минимальная ОИСКО}

Оптимальный параметр определяется из условия минимума:
\[
\xi_{\mathrm{opt}}(n)=\arg\min_{\xi>0}\bar{\delta}_n(\xi),
\qquad
\bar{\delta}_{n,\min}=\bar{\delta}_n(\xi_{\mathrm{opt}}).
\]
Для рассматриваемого значения $n=400$ минимум находится в первой ветви.
Дифференцируя первую ветвь по $\xi$, получаем
\[
\frac{d}{d\xi}\bar{\delta}_n(\xi)
=\frac{1}{12a}+\frac{1}{6an}
-\frac{1}{2an\xi^2}.
\]
Из условия равенства производной нулю следует
\[
\xi_{\mathrm{opt}}(n)=\sqrt{\frac{6}{n+2}}.
\]
При $n=400$
\[
\xi_{\mathrm{opt}}(400)=\sqrt{\frac{6}{402}}\approx 0.122,
\qquad
\bar{\delta}_{400,\min}\approx 0.0111.
\]

Эффективность параметра $\xi$ определим отношением
\[
e_n(\xi)=\frac{\bar{\delta}_{n,\min}}{\bar{\delta}_n(\xi)},
\qquad 0<e_n(\xi)\le 1.
\]
Максимальное значение эффективности равно $1$ и достигается при
$\xi=\xi_{\mathrm{opt}}(n)$.

\section*{Графическая иллюстрация}

Основные графики построены для $n=400$ по формуле
$\bar{\delta}_n(\xi)$ без дополнительных аппроксимаций.

\begin{figure}[H]
\centering
\plotplaceholder[0.82\textwidth]{plots/delta_bar_n400.png}
\caption{График ОИСКО $\bar{\delta}_{400}(\xi)$, построенный по аналитической формуле для прямоугольного ядра.}
\end{figure}

При малых $\xi$ ошибка быстро растёт из-за дисперсионного члена
$\sim (2an\xi)^{-1}$, а при больших $\xi$ кривая приближается к пределу
$1/(2\sqrt{3})$.

\begin{figure}[H]
\centering
\plotplaceholder[0.82\textwidth]{plots/efficiency_en_n400.png}
\caption{Эффективность $e_{400}(\xi)=\bar{\delta}_{400,\min}/\bar{\delta}_{400}(\xi)$ с максимумом, равным $1$.}
\end{figure}

Эффективность максимальна вблизи оптимального значения $\xi$ и убывает как
при слишком малых, так и при слишком больших значениях параметра.

\begin{figure}[H]
\centering
\plotplaceholder[0.82\textwidth]{plots/xi_opt_vs_n.png}
\caption{Дополнительно: зависимость $\xi_{\mathrm{opt}}(n)$ от размера выборки $n$.}
\end{figure}

По мере роста объёма выборки оптимальное значение $\xi$ убывает, поскольку
для больших $n$ можно использовать более узкое окно без сильного штрафа по
дисперсии.

\begin{figure}[H]
\centering
\plotplaceholder[0.82\textwidth]{plots/delta_min_vs_n.png}
\caption{Дополнительно: зависимость $\bar{\delta}_{n,\min}$ от размера выборки $n$.}
\end{figure}

Минимальная ОИСКО монотонно убывает: увеличение выборки улучшает наилучшее
достижимое качество оценки.
